% SHARD-ID: AQM-61-TANGENT-COTANGENT-WEYL-KINEMATICS-EXAMPLES
% SHARD-TITLE: Tangent-Cotangent Weyl Kinematics: Recipe and First Examples
% SHARD-SUMMARY: Defines the tangent-cotangent Weyl kinematical recipe with Hilbert space \(\ell^2(T_PX)\).
% SHARD-SUMMARY: Computes the first five finite \(\mathbb F_3\) examples with explicit tangent, cotangent, phase, and Hilbert spaces.
% SHARD-SUMMARY: Separates this kinematics from any de Rham CSS or Hamiltonian dynamics.
% SHARD-KEYWORDS: tangent, cotangent, Weyl, Hilbert space, F3, circle, singularity, nilpotent, kinematics

\section{Tangent-Cotangent Weyl Kinematics: Recipe and First Examples}
\label{sec:tangent-cotangent-weyl-kinematics-examples}

\subsection{Status and Source Anchors}

\begin{sloppypar}
This shard starts the kinematical recipe named in convention (ao) as
\(\texttt{tangent\_cotangent\_weyl\_kinematics}\).  K\"ahler differentials and
cotangent fibres are sourced from
\path{references/algebraic_geometry/stacks_project_algebra.tex}, lines
33793--33872 and 34310--34340.  Tangent spaces by derivations and dual numbers
are sourced from
\path{references/algebraic_geometry/stacks_project_varieties.tex}, lines
2915--3019.  Finite Weyl operators and additive characters are the AQM-28
convention, sourced there from Ashikhmin--Knill.  AQM-60 separates this
kinematics from the finite \(\texttt{de\_rham\_css\_test}\).
\end{sloppypar}

\subsection{Recipe}

Let \(X=\Spec A\) be a \(k\)-scheme presented by a \(k\)-algebra, and let
\(P\in X\) be a finite-residue point with
\[
  K=\kappa(P).
\]
The cotangent and tangent spaces are
\[
  T_P^*X=\Omega^1_{A/k}\otimes_AK,\qquad
  T_PX=\operatorname{Hom}_K(T_P^*X,K).
\]
Set
\[
  E_P^{\rm geom}=T_PX\oplus T_P^*X
\]
with the canonical symplectic form
\[
  \omega((v,\alpha),(w,\beta))=\beta(v)-\alpha(w).
\]
The Hilbert space is the configuration Hilbert space
\[
  \mathcal H_P^{\rm geom}=\ell^2(T_PX).
\]
For finite \(K\), use
\[
  \psi_K(a)=\exp\!\left(\frac{2\pi i}{p}
    \operatorname{Tr}_{K/\mathbb F_p}(a)\right).
\]
The Weyl generators act on the basis \(|v\rangle\), \(v\in T_PX\), by
\[
  T(w)|v\rangle=|v+w\rangle,\qquad
  R(\alpha)|v\rangle=\psi_K(\alpha(v))|v\rangle .
\]
Thus \(W(w,\alpha)=T(w)R(\alpha)\) is the finite Weyl-Heisenberg
prequantisation of the tangent-cotangent phase space.  If
\(\dim_KT_PX=m\), then
\[
  \mathcal H_P^{\rm geom}\cong \ell^2(K^m),\qquad
  \dim_{\mathbb C}\mathcal H_P^{\rm geom}=|K|^m.
\]
After choosing a \(K\)-basis of \(T_PX\), this is \(m\) qudits of local
dimension \(|K|\).

For a finite set \(S\) of finite-residue points, define
\[
  E_S^{\rm geom}=\bigoplus_{P\in S}E_P^{\rm geom},\qquad
  \mathcal H_S^{\rm geom}=\bigotimes_{P\in S}\mathcal H_P^{\rm geom}.
\]
This is kinematics only.  No stabilizer, Hamiltonian, or dynamics is selected
in this shard.

\subsection{Presentation Rule}

For a presentation
\[
  A=k[x_1,\ldots,x_n]/(f_1,\ldots,f_r)
\]
and a \(K\)-valued point \(P=(a_1,\ldots,a_n)\), a tangent vector is
\[
  v=(v_1,\ldots,v_n)\in K^n
\]
such that
\[
  \sum_i \frac{\partial f_j}{\partial x_i}(P)v_i=0
  \qquad\text{for every }j.
\]
The cotangent space is the dual quotient generated by \(dx_i\) modulo the
linear relations
\[
  \sum_i \frac{\partial f_j}{\partial x_i}(P)dx_i=0.
\]

\subsection{Examples 1--5}

\paragraph{1. The point \(\Spec\mathbb F_3\).}
Let \(A=k=\mathbb F_3\).  There is one point \(P\), \(K=k\), and
\[
  \Omega^1_{k/k}=0.
\]
Hence
\[
  T_P^*X=0,\qquad T_PX=0,\qquad E_P^{\rm geom}=0.
\]
The Hilbert space is
\[
  \mathcal H_P^{\rm geom}=\ell^2(0)\cong\mathbb C,
\]
with the single basis vector \(|0\rangle\).  There is no local qutrit in this
geometric kinematics.

\paragraph{2. The product \(k^2\).}
Let \(A=k^2\) with \(k=\mathbb F_3\).  AQM-58 proves
\[
  \Omega^1_{A/k}=0.
\]
There are two points \(P_1,P_2\), both with residue field \(K=k\), but at each
point
\[
  T_{P_i}^*X=0,\qquad T_{P_i}X=0.
\]
Thus
\[
  \mathcal H_{P_i}^{\rm geom}\cong\mathbb C,\qquad
  \mathcal H_{\{P_1,P_2\}}^{\rm geom}\cong\mathbb C.
\]
This differs from residue-Weyl kinematics, which would attach one
three-level residue qudit to each point.

\paragraph{3. Dual numbers over \(\mathbb F_3\).}
Let
\[
  A=k[\epsilon]/(\epsilon^2),\qquad k=\mathbb F_3,
\]
and let \(P=(\epsilon)\).  Since
\[
  d(\epsilon^2)=2\epsilon\,d\epsilon,
\]
the fibre at \(P\) kills \(\epsilon\), so
\[
  T_P^*X=k\,d\epsilon,\qquad T_PX=k\,\partial_\epsilon.
\]
Therefore
\[
  E_P^{\rm geom}=k\,\partial_\epsilon\oplus k\,d\epsilon,
  \qquad \omega((r,s),(r',s'))=s'r-sr'.
\]
The Hilbert space is
\[
  \mathcal H_P^{\rm geom}=\ell^2(k)
  =\operatorname{span}_{\mathbb C}\{|0\rangle,|1\rangle,|2\rangle\}
  \cong\mathbb C^3.
\]
The operators are \(T(a)|r\rangle=|r+a\rangle\) and
\(R(b)|r\rangle=\exp(2\pi i br/3)|r\rangle\).

\paragraph{4. The cubic fat point in characteristic three.}
Let
\[
  A=k[t]/(t^3),\qquad k=\mathbb F_3,\qquad P=(t).
\]
Here
\[
  d(t^3)=3t^2dt=0,
\]
so the relation contributes no first-differential relation.  Thus
\[
  T_P^*X=k\,dt,\qquad T_PX=k\,\partial_t.
\]
Again
\[
  \mathcal H_P^{\rm geom}=\ell^2(k)\cong\mathbb C^3.
\]
The length-three nilpotent thickening is invisible to the first tangent
dimension beyond saying that one tangent qutrit exists.

\paragraph{5. Two independent dual directions.}
Let
\[
  A=k[u,v]/(u^2,v^2),\qquad k=\mathbb F_3,\qquad P=(u,v).
\]
The relations give \(2u\,du=0\) and \(2v\,dv=0\).  At \(P\), both \(u\) and
\(v\) vanish, so
\[
  T_P^*X=k\,du\oplus k\,dv,\qquad T_PX=k\,\partial_u\oplus k\,\partial_v.
\]
The phase space is \(k^2\oplus(k^2)^*\), and
\[
  \mathcal H_P^{\rm geom}=\ell^2(k^2)\cong\mathbb C^9
  \cong\mathbb C^3\otimes\mathbb C^3.
\]
A basis vector is \(|r,s\rangle\) with \((r,s)\in\mathbb F_3^2\).

Examples 6--10 continue in AQM-62, including affine-line, affine-plane,
crossing, and finite-field circle supports.
